\clearpage
\VerdierPhysicalPageAnchor{0247}%
\begingroup\small\setlength{\parskip}{0pt}\setlength{\emergencystretch}{1.5em}
\setlength{\abovedisplayskip}{4pt plus 2pt minus 2pt}%
\setlength{\belowdisplayskip}{4pt plus 2pt minus 2pt}%

\VerdierSemanticAnchor{ch4d-4-6-5-resolution-type}
\noindent soient des résolutions de type injectif (1.1.4).

\VerdierSemanticAnchor{ch4d-4-6-5-terminology}
Nous appelons donc \emph{résolution injective de Cartan--Eilenberg du
complexe $Y$} ce que ces auteurs appellent simplement \emph{résolution
injective du complexe $Y$} ([1], chap.~XVII), terminologie que nous
n'utiliserons pas dans ce travail. Une résolution injective de
Cartan--Eilenberg de $Y$ induit des résolutions de type injectif des
composants du complexe $Y$ et des noyaux des différentielles de $Y$
[\emph{loc.\ cit.}].

\VerdierSemanticAnchor{ch4d-4-6-5-finite-length}
Nous dirons qu'une résolution de Cartan--Eilenberg
$\varepsilon\colon Y\longrightarrow I^{\bullet\bullet}$ est de
\emph{longueur finie} s'il existe un entier $n_0$ tel que
$I^{\bullet,n}=0$ pour $n>n_0$. Pour qu'un complexe $Y$ admette une
résolution injective de Cartan--Eilenberg de longueur finie, il faut et il
suffit que la dimension injective des objets de cohomologie de $Y$ et des
images des différentielles de $Y$ soit majorée. Rappelons les résultats
établis dans \emph{loc.\ cit.} :

\VerdierSemanticAnchor{ch4d-4-6-6-a}
\medskip\noindent\textbf{Proposition 4.6.6.}\quad
\textit{a) Lorsque la catégorie $\mathcal A$ possède suffisamment
d'objets injectifs (3.1.5), tout complexe simple $Y$ de $\mathcal A$ admet
une résolution injective de Cartan--Eilenberg.}

\VerdierSemanticAnchor{ch4d-4-6-6-b}
\noindent\textit{b) Soient
$f\colon X\longrightarrow Y$ un morphisme de complexes simples,
$\varepsilon\colon X\longrightarrow I^{\bullet\bullet}$ et
$\varepsilon'\colon Y\longrightarrow J^{\bullet\bullet}$ deux résolutions
injectives de Cartan--Eilenberg. Il existe un morphisme de doubles
complexes :}

\VerdierSemanticAnchor{ch4d-4-6-6-b-morphism}
\[
\textit{$g\colon I^{\bullet\bullet}\longrightarrow J^{\bullet\bullet}$}
\]
\textit{au-dessus de $f$, i.e.\ tel que le diagramme de morphismes de
doubles complexes :}

\VerdierSemanticAnchor{ch4d-4-6-6-b-diagram}
\[
\begin{tikzcd}[column sep=5.0em,row sep=3.0em]
X \arrow[r,"f"] \arrow[d,"\varepsilon"'] &
Y \arrow[d,"\varepsilon'"] \\
I^{\bullet\bullet} \arrow[r,"g"'] &
J^{\bullet\bullet}
\end{tikzcd}
\]
\textit{soit commutatif.}

\VerdierSemanticAnchor{ch4d-4-6-6-c}
\noindent\textit{c) Deux morphismes
$g,g'\colon I^{\bullet\bullet}\longrightarrow J^{\bullet\bullet}$
au-dessus de $f$ sont homotopes (au sens de l'homotopie des morphismes de
doubles complexes (chap.~I, 2.5.1)).}

\VerdierSemanticAnchor{ch4d-4-6-6-proof}
Nous renvoyons pour la démonstration à \emph{loc.\ cit.}

\VerdierSemanticAnchor{ch4d-4-6-7}
\medskip\noindent\textbf{4.6.7.}\quad
Soient $Y$ un objet de $\operatorname{comp}(\mathcal A)$ et
$Y\longrightarrow I^{\bullet\bullet}$ une résolution injective de
Cartan--Eilenberg de $Y$. Le double complexe $I^{\bullet\bullet}$ donne
naissance à deux objets spectraux à valeurs dans $D(\mathcal A)$, suivant
qu'on le filtre par son premier ou son deuxième degré (4.1.7). L'objet
spectral, à valeurs dans $D(\mathcal A)$, défini par
$I^{\bullet\bullet}$ filtré par son deuxième degré, est appelé
\emph{l'objet spectral de Cartan--Eilenberg de $Y$} (relatif à la
résolution de Cartan--Eilenberg considérée).
\endgroup
